\section{The experimental problem}\label{sec:problem}
A causal experiment does not present an observer with a parameter manifold. It presents preparations, admissible interventions and reports. Geometry becomes relevant only after specifying which future reports are to be predicted and which information the observer has actually acquired. There is a further obstruction: a small terminal codebook need not be recursively selectable after earlier reports have been discarded. We use \emph{General Theta} for the mathematical program connecting prepared causal kernels, executable future tests, their minimal predictive quotient, resource-certified experiment comparison, acquired geometry and prediction or decision risk. ``Foundations'' refers to this order of construction, not to a claim that all causal encoders have already been classified.

The central result is Theorem~\ref{thm:serial}. Suppose the acquired response at cut $t$ has a positive-mass observable submeasure of exponent $s_t$, and a globally selectable cover of the same order is stable under the raw task updates. The stability requirement includes a Lyapunov envelope for all candidates, not just the exact trajectory or the code eventually chosen. Then, for one fixed terminal task and its full-history Bayes risk $B_T$,
\[
 R_\on(T,M)-B_T\asymp M^{-2/s_*},\qquad
 R_\cp(T,M)-B_T\asymp M^{-2/s_T},\qquad s_*=\max_t s_t.
\]
The entire family of actual continuation widths is equivalent in order to the online risk. Within this declared class, checkpoint and online compression match if and only if $s_*=s_T$. The theorem allows noninteger exponents, changing state dimension, dependent reports and expansive finite-horizon updates. Its proof constructs one $M$-label recursion at every cut; it does not assume that separately optimal checkpoint codebooks can be combined.

Two mechanisms make this possible. A conditional future response is a function of the information arriving after a cut. An earlier label must therefore choose a function, not merely a terminal scalar. Quantizing the actual dominated measure of these functions yields a common-task lower. Theorem~\ref{thm:cutduality} identifies the exact one-cut functional-coding problem, compares it with acquired Hilbert measures, and localizes it when physically shared side information prevents global product domination. Raw density and joint regeneration lemmas produce its common continuation laws. On the constructive side, global covers and a candidate Lyapunov bound control every inserted error under the actual report law, without independence or a free exact posterior register. The block transfer theorem, Theorem~\ref{thm:transfer}, retains its separate strength of time-uniform control under verified raw-block stability.

Theorem~\ref{thm:ranknoisy} verifies the central theorem from positive-noise correlated Gaussian kernels. Its successive observable ranks $r_j$ can change, and its largest rank $r$ is determined by the executable query family. It proves online excess $M^{-2/r}$ against checkpoint excess $M^{-2}$, even in an ambient hidden dimension $d>r$. A partial task statistic is not mislabeled a full predictive quotient: nuisance directions can still affect the suffix law. Theorem~\ref{thm:singularserial} instead starts from a Cantor-product preparation of total exponent $D$, obtained from its actual cylinder masses. A noisy continuation and continuous terminal query produce online excess $M^{-2/\max(D,1)}$. The terminal cut is essential when $D<1$. These two realizations are consequences of the same general theorem, not premises from which it is inferred.

Acquisition itself is a different optimization. A fixed exploration lower cannot be optimized by interchanging infima over its policy-dependent acquired law. Theorem~\ref{thm:adaptive} resolves this problem for progressive sensors with independent fair unread tails and arbitrary nonhomogeneous scales. The observed bit values do not alter future tail scales; this is why deterministic greedy allocation attains the optimum over all admitted adaptive allocations. Theorem~\ref{thm:joint} matches acquisition, retained memory, calibration and an attained simulation defect in one task. Theorem~\ref{thm:digital} gives a finite-bit feasible account, not a general computational lower. Theorem~\ref{thm:brier} preserves the sharp persistent-state law for ordinary filtering loss in sparse hidden-state models. A sequential-chart theorem and a full-rank noisy regression theorem are retained as intermediate consequences and independently checkable specializations.

The compatibility criterion requires both acquired observability and globally implementable update certificates. It does not derive the latter from bare one-cut widths for every causal experiment. The matched acquisition theorem does not solve unknown-kernel or arbitrary observation-dependent experimental design. Conditional projection, functional source coding, quantization, classical allocation and experiment comparison are explicitly credited in Section~\ref{sec:literature}. Supplement S~\cite{companion} is part of this same submission, supplied in full; it preserves the nonreset singular, stopped, calibration, output-grid and typed-interface proofs. It is not an external publication or a repository-only proof premise.

\subsection{Kernels, interfaces and resources}
All measurable spaces are standard Borel. An experiment is specified by a preparation $\mu_\lambda$, hidden states $X_t$, actions $A_t$, reports $Y_{t+1}$, and jointly measurable probability kernels
\begin{equation}\label{eq:kernel}
 K_t^\lambda(dx',dy,dc\mid x,a).
\end{equation}
Here $c$ records nonnegative raw-call, time and other declared resource increments. Positivity means nonnegativity and normalization. Dirac operations, failure symbols and missing reports are allowed. A visible history $h_t$ consists of the preparation report, public exploration coins, and intervening actions, reports and public costs. Private exploration coins that influence later actions remain in the physical/controller state and the actual conditional law unless genuinely reported. Legal action sets have a declared measurable graph. Strategies are parameter-independent Borel kernels on those sets; stopping rules use the visible filtration. They cannot consult future reports or the unknown $\lambda$.

A continuation interface is a measurable function $v_t(h_t)$ which specifies legal future actions, available resources and any publicly supplied phase. A predictive realization is a standard Borel state $S_t=s_t(h_t)$ determining that interface, every declared future-test expectation, and jointly Borel report and update kernels $J_t,F_t$. When a phase is kept separate, write the state space as a measurable disjoint union $S=\coprod_{v\in V}S_v$ over a \emph{finite} public phase set. All comparison pairs at a given time lie in the same fiber $S_v$. Updates map both members into the same next fiber, whose phase is a declared function of the old phase, action and common report. No unretained data-dependent phase is public for free. An unbounded or data-dependent external interface requires a separately stated conditional theorem and is not covered by the alphabet count below.

A predictor retains at most $M$ data labels at each declared cut. Its update uses only its current retained tuple, current action/report and allowed fresh randomness. Any controller, simulator or clock register consulted later is counted separately or included in the label tuple. Temporary workspace must be erased before the next cut; program constants must be known, parameter-independent data with a charged description. A vector output is a write-only stream: previously emitted coordinates are not rereadable workspace. The notation $M$ is an alphabet cardinality, not total arithmetic space. If a public finite phase has $Q$ states, the full serial cardinality is at most $MQ$; a lower bound at a cut uses the \emph{entire} accessible data-dependent tuple. Deterministic time is fixed in the main prediction statements and is not a hidden observation channel.

For the general transfer theorem fix a preparation and a legal exploration independently of the prediction encoder. Algorithm coins are independent of that experiment. In the adaptive acquisition theorem the exploration is instead optimized together with the encoder; its lower bounds are proved anew and do not condition an encoder-dependent law as though it were fixed.

\subsection{Executable tests and a measurable predictive quotient}
An executable test is a legal finite continuation followed by a bounded function of its scored reports. Two histories are equivalent if their continuation interfaces and all test expectations agree. A set-theoretic quotient need not be a suitable measurable state. We record a sufficient realization result.

\begin{lemma}[A compact measurable section]\label{lem:selector}
A continuous surjection $e:K\to S$ between compact metrizable spaces has a Borel section.
\end{lemma}
\begin{proof}
For each $n$ fix a finite closed cover $(F_{n,j})_j$ of $K$ by nonempty sets of diameter at most $2^{-n}$. Inductively select the first $j_n$ for which $e^{-1}(s)\cap F_{1,j_1}\cap\cdots\cap F_{n,j_n}$ is nonempty. For a fixed finite index word its compact intersection has closed image under $e$. Each finite-word selection event is therefore Borel, using finite Boolean combinations of those closed images. The selected compact intersections are nested, meet the fiber, and have diameter tending to zero. They have a unique common point in $e^{-1}(s)$. Choosing once a point of each nonempty finite-word intersection gives Borel finite-valued approximants converging to that point. Their pointwise limit is a Borel section.
\end{proof}

\begin{proposition}[Continuous-instrument realization]\label{prop:quotient}
Let $\mathcal B$ be compact metrizable and let $\mathcal V\subset C(\mathcal B)$ be the separable closed linear span of constants and the continuous evaluations of a determining family of executable tests and interface coordinates. All kernel, density, update and instrument versions below are jointly Borel also in the action variable on the declared standard Borel action--report domain. For each fixed action, let its report space be Polish, with reference measure $m_a$, and let $g_a(\pi,y)$ be a jointly continuous density. On the open set $D_a=\{(\pi,y):g_a(\pi,y)>0\}$, with the relative product topology, suppose the normalized update $B_a$ is continuous. For every $f\in\mathcal V$, require the instrument evaluation $g_a(\pi,y)f(B_a(\pi,y))$, with its specified value at density zero, to belong to $\mathcal V$ as a function of $\pi$ and to be jointly continuous. Then the countable evaluation image is a compact metrizable predictive quotient. Report laws and updates descend, and the latter are continuous on the descended open positive-density domain. Every standard Borel statistic measurably determining all these evaluations measurably determines the quotient.
\end{proposition}
\begin{proof}
Normalize a countable dense family of evaluations and map $\mathcal B$ into $[-1,1]^{\N}$. The image $S$ is compact, hence Borel. Equality of its coordinates gives equality on $\mathcal V$ by uniform approximation and hence on the determining tests. Conversely, equality of all future tests and interface coordinates gives equality on their closed linear span. Thus the fibers are exactly predictive equivalence classes, not a finer partition carrying unrelated coordinates. Taking $f=1$ shows that density values agree on each fiber. Dividing equal instrument values by that positive common density shows that the evaluation of the updated belief agrees on the fiber.

The evaluation map is a continuous surjection from a compact space to a Hausdorff space, hence is perfect. Its product with the report identity is also a quotient map: closedness follows by taking a convergent subnet in the compact belief domain whenever an image net converges. The instrument coordinates and their positive-density ratios therefore descend continuously. The positive-density domain is open in $S\times Y$; no update is defined by division at a null report. For measurability jointly across actions, Lemma~\ref{lem:selector} supplies a Borel section of the compact evaluation map. Composing the assumed joint versions with this section supplies joint Borel versions of all descended kernels. A common jointly Borel null-report extension, when needed on wrong candidates, is additional data.

For minimality, collect the measurable coordinate factors through the given statistic into a measurable product map. It takes values in the Borel set $S$ on the domain in question; extend outside that set by a fixed point. Countably many almost-sure factors can be imposed on a single common full-measure history set. This proves measurable, rather than merely set-theoretic, minimality.
\end{proof}

Applications may verify a predictive realization directly. The metric used for a quantitative estimate need only induce its standard Borel structure; it need not be the compact topology in Proposition~\ref{prop:quotient}.

\subsection{One prediction task}
Let $Y$ be a bounded terminal vector in a finite-dimensional Euclidean space with a specified weighted norm. A finite independent terminal probe index is interpreted as the vector of its separately executable conditional means, with the probe probabilities as weights; it is not a simultaneous counterfactual measurement. At the decision cut $T$ put
\begin{equation}\label{eq:task}
 P_T=\E[Y\mid\mathcal H_T],\qquad B_T=\E\norm{Y-P_T}^2,
 \qquad \nu_T=\law(P_T).
\end{equation}
For a finite Borel measure $\nu$ on a separable Hilbert space define
\begin{equation}\label{eq:quantization}
 q_M(\nu)=\inf_{c_1,\ldots,c_M}\int\min_i\norm{x-c_i}^2\,\nu(dx).
\end{equation}
The measure may be atomic, singular or a subprobability. A checkpoint algorithm inspects $\mathcal H_T$ once before retaining $M$ labels. An online algorithm retains at most $M$ labels at every earlier cut and cannot reread a report. In both cases the terminal outcome is revealed only after the prediction. Denote the optimal risks by $R_\cp(T,M)$ and $R_\on(T,M)$.
